\begin{center}
	\begin{longtable}{|p{5.1cm}|p{5.1cm}|p{5.1cm}|}
		\hline
		\textbf{Pouvez-vous citer trois points positifs du plugin ? } & \textbf{Pouvez-vous citer trois points negatifs du plugin ? } & \textbf{Quels succès auriez-vous souhaité voir dans le plugin ?} \\
		\hline
		Facile, intuitive, intéressant & setup initial peut être complexe pour le débutants,   & nan \\
		\hline
		Intégré, Dynamique,  & Manque d'explication pour les goals, Goals devraient varier par classes,  Permettre un suivi en live & nan \\
		\hline
		Drôle, motivant, curieux & Une fois les succès réalisé, l'intérêt disparait, Ne fonctionne que avec le code fourni (classe dragon), bugs de mise à jour & nan \\
		\hline
		icone sympathiques, facile à utiliser, donne envie de compléter & pas beaucoup de succès & nan \\
		\hline
		ludique, entrainant, facile d'accès & evaluation nombre de test -> écriture de test inutiles, idem pour coverage, peut être difficile à adapter sur un projet réel & nan \\
		\hline
		débloquer des achievements est chouette, leurs noms sont amusants, c'est simple à utiliser & C'est marrant au début mais je pense qu'à long terme on peut se lasser, j'aimerais avoir plus d'achievements & nan \\
		\hline
		Facile à utiliser, facile à installer & Incompatibilité avec mes versions IDE et jdk d'avant & nan \\
		\hline
		nouveaux , sympa , constructif  & complexe à prendre en main , perdu , pas simple  & aucune idée \\
		\hline
		Les objectifs clairs, l'excellente intégration à IntelliJ (si on oublie les bugs), les objectifs personnalisés au projet (ex: nettoyer les dragons de leurs mutants, etc.) & difficilement adaptable au projet, une partie des succès peut-être réalisée avec des tests poubelle et inutiles, plus gadget qu'autre chose dans un vrai projet & Rien de particulier en plus de ce qu'il y a déjà \\
		\hline
		UX, User friendly et bien documenté & les bugs & un pour ce combatre au sein de la team \\
		\hline
		amusant, stimulant et claire (j'aime les petits logo des achievements) & bug avec les mutants; mode equipe moins stimulant( j'aurais preferer le mode solo) ; fluidité & nombre de teste reussie avec plus de 90 % de coverage ; \\
		\hline
		facile en prendre en main, UX propre, et succes unlcok rapidement & dependance des autres membres d'equipes, sentiment de competition qui peut distraire, et quelques bugs  & success pour l'accomplissement de groupe en synergie avec tout le membre (quand un certain palier est atteint) \\
		\hline
		Bien intégré, sobre, les succès sont motivant & Manque de transparence sur le moment de la mise à jour, les images des succès n'était pas toujours claires, la description des succès n'étaient pas toujours claires & Un succès sur le coverage global du projet \\
		\hline
	\end{longtable}
	\label{table:open-questions-answers}
\end{center}